% ======================================================================
% Section 3.5 -- Database Design (ERD)
% ======================================================================

\section{Entity Identification}
\begin{table}[H]
\centering
\caption{Entity Identification}
\label{tab:entities}
\begin{tabularx}{\textwidth}{@{}lX@{}}
\toprule
Entity & Description \\
\midrule
\textbf{User} & System accounts --- UP officials, Upazila officers, NGO workers \\
\textbf{GeographicArea} & Full BD administrative hierarchy: Division $\rightarrow$ District $\rightarrow$ Upazila $\rightarrow$ Union $\rightarrow$ Ward \\
\textbf{Household} & Registered disaster-affected family unit with age-bracket counts \\
\textbf{DistributionLog} & Record of one relief distribution (optionally linked to a pledge) \\
\textbf{ItemCategory} & Lookup table for predefined relief item types with per-person rates \\
\textbf{DuplicateAlert} & Alert generated when duplicate distribution is detected \\
\textbf{SyncConflict} & Log of offline sync conflicts pending manual review \\
\textbf{Feedback} & User-submitted feedback/complaints with admin response \\
\textbf{Inventory} & Stock tracking per item category \\
\textbf{NeedAssessment} & Ward-level calculated need per item category, with optional officer override \\
\textbf{ReliefPledge} & Source declaration of relief supply with status lifecycle \\
\bottomrule
\end{tabularx}
\end{table}

\section{Entity-Relationship Diagram}
The ERD uses Crow's Foot notation with the relationships shown below:

\begin{center}
\begin{tikzpicture}[
    ent/.style={rectangle, draw=rmstublue3!80!black, thick, fill=rmstublue3!5, minimum width=2.6cm, minimum height=0.7cm, align=center, font=\small\bfseries, rounded corners=3pt},
    rel/.style={-{Stealth[length=2mm,width=2mm]}, draw=rmstugray, thick},
    lab/.style={font=\tiny, fill=white, inner sep=1pt}
]
% Entities
\node[ent] (GA) {GeographicArea};
\node[ent, below=0.8cm of GA] (U) {User};
\node[ent, below left=1.8cm and 2.5cm of GA] (HH) {Household};
\node[ent, below right=1.8cm and 2.5cm of GA] (IP) {ItemCategory};
\node[ent, below=1.5cm of HH] (DL) {DistributionLog};
\node[ent, below right=0.8cm and 1.5cm of DL] (RP) {ReliefPledge};
\node[ent, below left=0.8cm and 1.5cm of DL] (NA) {NeedAssessment};

% Relationships
\draw[rel] (GA) -- (U) node[midway,left,lab] {1..M};
\draw[rel] (GA) -- (HH) node[midway,above,sloped,lab] {1..M};
\draw[rel] (GA) -- (NA) node[midway,left,lab,font=\tiny] {1..M};
\draw[rel] (GA) -- (RP) node[midway,right,lab,font=\tiny] {1..M};
\draw[rel] (HH) -- (DL) node[midway,right,lab] {1..M};
\draw[rel] (U) -- (DL) node[midway,above,sloped,lab] {1..M};
\draw[rel] (IP) -- (DL) node[midway,above,sloped,lab,font=\tiny] {1..M};
\draw[rel] (IP) -- (NA) node[midway,right,lab,font=\tiny] {1..M};
\draw[rel] (IP) -- (RP) node[midway,above,sloped,lab,font=\tiny] {1..M};
\draw[rel, dashed] (DL) -- (RP) node[midway,above,lab] {M..1};
\end{tikzpicture}
\captionof{figure}{Entity-Relationship Diagram (simplified)}
\end{center}

\paragraph{Key Relationships}

\begin{itemize}[nosep]
    \item \textbf{GeographicArea} (1) $\longleftrightarrow$ (Many) \textbf{User} --- one area has many users
    \item \textbf{GeographicArea} (1) $\longleftrightarrow$ (Many) \textbf{Household} --- one union has many households
    \item \textbf{GeographicArea} (self-referencing) --- parent\_id for hierarchy
    \item \textbf{Household} (1) $\longleftrightarrow$ (Many) \textbf{DistributionLog} --- one household receives many distributions
    \item \textbf{User} (1) $\longleftrightarrow$ (Many) \textbf{DistributionLog} --- one officer logs many distributions
    \item \textbf{DistributionLog} (Many) $\longleftrightarrow$ (1) \textbf{ReliefPledge} --- optional FK for pledge fulfillment
    \item \textbf{ItemCategory} (1) $\longleftrightarrow$ (Many) \textbf{DistributionLog}, \textbf{NeedAssessment}, \textbf{ReliefPledge}
    \item \textbf{GeographicArea} (1) $\longleftrightarrow$ (Many) \textbf{NeedAssessment}, \textbf{ReliefPledge}
\end{itemize}

\section{Normalization}
All tables are in \textbf{3NF (Third Normal Form)}:

\begin{itemize}[nosep]
    \item \textbf{1NF:} All attributes atomic --- vulnerability flags are separate boolean columns, age-bracket counts are separate INT columns
    \item \textbf{2NF:} All tables use single-column surrogate PKs (UUID); no composite PKs
    \item \textbf{3NF:} GeographicArea hierarchy is self-referencing; ItemCategory is in its own table; ReliefPledge.remaining\_qty is a computed column
\end{itemize}

\section{Data Dictionary (Key Tables)}

\subsection{users}
\begin{itemize}[nosep]
    \item user\_id (PK), area\_id (FK), name, email (UNIQUE), password\_hash, role (ENUM), organization, is\_active, created\_at
\end{itemize}

\subsection{geographic\_areas}
\begin{itemize}[nosep]
    \item area\_id (PK), name, level (ENUM: DIVISION/DISTRICT/UPAZILA/UNION/WARD), parent\_id (FK, self), geometry (GeoJSON)
\end{itemize}

\subsection{households}
\begin{itemize}[nosep]
    \item hh\_id (PK), area\_id (FK), head\_name, nid (UNIQUE), gps\_lat, gps\_lng, family\_size, is\_elderly, is\_disabled, is\_pregnant, children\_0\_5, over\_60, photo\_url, registered\_by (FK), created\_at
\end{itemize}

\subsection{distribution\_logs}
\begin{itemize}[nosep]
    \item log\_id (PK), hh\_id (FK), officer\_id (FK), item\_category\_id (FK), pledge\_id (FK, nullable), quantity, unit, gps\_lat, gps\_lng, photo\_url, timestamp, sync\_status (ENUM), is\_override, override\_reason
\end{itemize}

\subsection{need\_assessments}
\begin{itemize}[nosep]
    \item assessment\_id (PK), area\_id (FK, level=WARD), item\_category\_id (FK), calculated\_qty, override\_qty (nullable), override\_reason, computed\_at, overridden\_at, overridden\_by (FK)
\end{itemize}

\subsection{relief\_pledges}
\begin{itemize}[nosep]
    \item pledge\_id (PK), source\_id, source\_type (ENUM), area\_id (FK), item\_category\_id (FK), pledged\_qty, remaining\_qty, team\_count, volunteer\_count, description, status (ENUM: PENDING/IN\_FULFILLMENT/COMPLETED/CANCELLED)
\end{itemize}
